\question{10}{Considere un grafo representado por la siguiente matriz de adyacencia:

\begin{center}
\(
\begin{pmatrix}
0 & 1 & 0 & 0 & 1 & 0 & 0 \\
1 & 0 & 1 & 0 & 0 & 1 & 0 \\
0 & 1 & 0 & 0 & 0 & 0 & 1 \\
0 & 0 & 0 & 0 & 1 & 0 & 1 \\
1 & 0 & 0 & 1 & 0 & 1 & 0 \\
0 & 1 & 0 & 0 & 1 & 0 & 0 \\
0 & 0 & 1 & 1 & 0 & 0 & 0
\end{pmatrix}
\)
\end{center}

La primera fila representa el vértice \(v_1\) y la última \(v_7\). Cada vez que visite los vecinos de un vértice, hágalo en orden ascendente de índice.
}

\begin{enumerate}
  \item (2 décimas) Dibuje el grafo que representa la matriz.
  \item (1 décimas) ¿El grafo es dirigido? ¿Por qué?
  \item (1 décimas) ¿El grafo es ponderado? ¿Por qué?
  \item (6 décimas) Ejecute DFS recursivo desde \(v_3\). Diligencie la tabla a continuación con el estado de las estructuras en el instante justo después de que se invoca un llamado al procedimiento recursivo de DFS. Liste los llamados a la función en la pila de llamados en el orden en el que ingresaron a la pila, de izquierda a derecha.

\begin{center}
\begin{tabularx}{\textwidth}{|>{\centering\arraybackslash}X|>{\centering\arraybackslash}X|>{\centering\arraybackslash}X|}
\hline
\textbf{Visitado} & \textbf{Pila de llamados} & \textbf{Predecesores} \\
\hline
$v_3$ \rule{0pt}{0.9cm} & & \\
\hline
\rule{0pt}{0.9cm} & & \\
\hline
\rule{0pt}{0.9cm} & & \\
\hline
\rule{0pt}{0.9cm} & & \\
\hline
\rule{0pt}{0.9cm} & & \\
\hline
\rule{0pt}{0.9cm} & & \\
\hline
\rule{0pt}{0.9cm} & & \\
\hline
\end{tabularx}
\end{center}
\end{enumerate}

\bigskip
